\documentclass[10pt,a4paper]{amsart}
\usepackage[margin=19mm]{geometry}
\usepackage{amsmath,amssymb,mathtools}
\usepackage[scr=rsfs,cal=euler]{mathalfa}
\usepackage{xfrac}
\usepackage[unicode,hidelinks,bookmarks=false]{hyperref}
\newcommand{\ie}{i.e.\ }
\newcommand{\cf}{cf.\ }
\newcommand{\resp}{respectively\ }
\newcommand{\Subsection}{\subsection}
\providecommand{\nobreakdash}{\nobreak\hbox{-}}
\setlength{\emergencystretch}{2em}
\multlinegap0pt
\usepackage{tikz}
\usetikzlibrary{cd}
\usepackage{mathtools}
\usepackage{amssymb}
\usepackage{xcolor}
\usepackage{tcolorbox}
\usepackage{enumerate}
\usepackage{tikz}
\newtheorem{theorem}{Theorem}
\newcommand{\Mod}{{\mathbf{Mod}}}
\newcommand{\Set}{\mathbf{Set}}
\newcommand{\T}{\mathbb{T}}
\newcommand{\colim}{\operatorname{colim}}
\newcommand{\id}{\operatorname{id}}
\title{De Morgan's law in toposes I}
\author{Olivia Caramello, Yorgo Chamoun}
\date{Source: arXiv:2507.01869v2 (CC BY 4.0)}
\begin{document}
\maketitle

\noindent\emph{Source and licence.} De Morgan's law in toposes I. Version arXiv:2507.01869v2 (CC BY 4.0), distributed by arXiv under the Creative Commons Attribution 4.0 International licence, http://creativecommons.org/licenses/by/4.0/, as recorded in the arXiv licence metadata for this exact version and shown on the abstract page https://arxiv.org/abs/2507.01869v2. Full licence notice: \texttt{licenses/arxiv-2507.01869-CC-BY-4.0.txt}.\par\smallskip

\noindent\textit{Editorial scope.} Counted unit: the article's theorem theorem (source0.tex lines 896--898). The article's complete original proof of it is delivered with it (source0.tex lines 899--933, one \texttt{proof} environment of 35 lines). No mathematics is written, repaired or re-derived: the statement and the proof are the article's own lines, in the article's own notation, and no retained line is substituted or re-flowed.  No further source-local context is needed: every cross-reference of every retained line resolves inside the two delivered spans.  Nothing was omitted from any retained span, so the package carries no omission record.  No retained line contains a citation, so the wrapper carries no bibliography and no bibliography entry.  The retained lines contain the article's own diagram code, which the wrapper typesets with the standard tikz packages; no image file is delivered and the package declares no asset dependency.  Editorial formatting only, in addition: an AMS \texttt{amsart} preamble that restates the article's own declarations and macros used by the retained lines, namely the environments \texttt{theorem} on separate counters, together with the standard packages those lines need.  The retained environments print in this document as Theorem 1 in order of appearance, whereas the article prints the same items with the numbering of its own full document.  The complete native source of the article as distributed in the arXiv e-print for this exact version is bundled unchanged as \texttt{sources/180910/source0.tex}.
\begin{theorem}
    Let $\T'\subseteq\T$ geometric theories of presheaf type. Suppose that the finitely presentable models of $\T$ are finitely presentable as models of $\T'$. Then $\T$ is De Morgan if and only if it has the weak relative amalgamation property with respect to $\T'$.
\end{theorem}

\begin{proof}
    Suppose that $f.p.\T$-$\Mod(\Set)$ has the amalgamation property. Consider $A\models\T$. Recall that we can write $A=\colim_{i\in I}G_i$ with $G_i$ finitely presented model of $\T$. By assumption, the $G_i$ are also finitely presented models of $\T'$. Now consider a diagram  
    \begin{center}
            \begin{tikzcd}
        G_i\arrow[r,"\id"] & G_i \arrow[r,""]\arrow[d] & H_1 \\
        & H_2 & 
    \end{tikzcd}
    \end{center}
    with $H_1,\, H_2\in FP(\T',\T)$, i.e. there is $h_1:H_1\xrightarrow{}A_1$ and $h_2:H_2\xrightarrow{}A_2$ with $A_1,\,A_2\models\T$. $h_i$ factors through some finitely presented model $H'_i$ of $\T$, and we get a diagram 
     \begin{center}
            \begin{tikzcd}
        G_i\arrow[r,"\id"] & G_i \arrow[r,""]\arrow[d] & H'_1 \\
        & H'_2 & 
    \end{tikzcd}
    \end{center}
    by composition, and this diagram admits amalgamation by hypothesis because $G,\,H'_1,\,H'_2$ are all finitely presented models of $\T$.\\
    \\
    Conversely, consider a diagram 
     \begin{center}
            \begin{tikzcd}
        H \arrow[r,"h_1"]\arrow[d, "h_2"'] & H_1 \\
        H_2 & 
    \end{tikzcd}
    \end{center}
    in $f.p.\T$-$\Mod(\Set)$. Write $H=\colim_{i\in I}G_i$ as given by hypothesis. The identity morphism $H\xrightarrow{}H$ factors through some $G_i$:
    $$\id=H\xrightarrow{f}G_i\xrightarrow{h}H$$
    and there is $g:G_i\xrightarrow{}G$ through which $h$ factors and we can amalgamate
    \begin{center}
            \begin{tikzcd}
        G_i\arrow[r,"g"] & G \arrow[r,""] & H \arrow[r,""]\arrow[d] & H_1 \arrow[d, dotted] \\
       & & H_2\arrow[r, dotted] & B
    \end{tikzcd}
    \end{center}
    We can then precompose with $f$ to get $\id_H$ before the square. This easily implies that we can amalgamate the diagram by a finitely presented model.
\end{proof}

\end{document}
